\documentclass[12pt]{exam}
\usepackage[T1]{fontenc}
\usepackage[utf8]{inputenc}
\usepackage[lithuanian]{babel}
\usepackage{lmodern}
\usepackage{graphicx}
\usepackage[a4paper, left=1cm, right=1cm, top=2.5cm, bottom=1cm]{geometry}
\usepackage{amsmath, amssymb}
\usepackage{tasks}

% Antraštė
\pagestyle{headandfoot}
\runningheadrule
\firstpageheader{Vardas: \rule{4cm}{0.4pt}}{Kontrolinis darbas: \\Daugianariai}{Vilties variantas}
\runningheader{Daugianariai}{}{}
\firstpagefooter{}{}{}
\runningfooter{}{}{}
\pointsdroppedatright          
\bracketedpoints               
\setlength{\rightpointsmargin}{2.5cm} 
\renewcommand{\points}{taškai} 
\renewcommand{\pointname}{Taškai:} 
\marginpointname{\ taškai}     
\renewcommand{\dottedlinefillheight}{0.7cm}
\newcommand{\atsakymas}{%
    \par\vspace{0.3cm}
    \textbf{Atsakymas:} \framebox[7cm]{\rule{0pt}{0.7cm}}
}
\renewcommand{\partshook}{%
    \setlength{\leftmargin}{1pt}%
    \setlength{\itemindent}{17pt}%
    \setlength{\labelsep}{0.5em}%
}

\begin{document}

\begin{questions}

\question[6] Nustatykite kiekvieno reiškinio tipą: racionalusis (R), trupmeninis (T), sveikasis (S) ar iracionalusis (I). Prie kiekvieno reiškinio nurodykite visas jam tinkančias raides.
\droppoints
\begin{tasks}[label-width=1.2em](3)
    \task $25^{-7}+x^2-x+\sqrt{6}$
    \task $\frac{3}{x ^{2} }-4x+7$
    \task $\sqrt{5-x}-2x+1$
    \task $\frac{x ^{2} }{x-1}+\sqrt{\pi}$
    \task $7x^{-2}-11$
    \task $\frac{2x+1}{\sqrt{11}}+x$
\end{tasks}



\question[5] Suprastinkite reiškinį ir nurodykite, ar gautas reiškinys yra vienanaris. Jei reiškinys yra vienanaris, nurodykite jo koeficientą ir laipsnį; jei ne, paaiškinkite kodėl.

\droppoints
    $$\displaystyle \frac{(-3x^7y^{-5})^{12}\cdot(27x^{-1}y^4)^{-20}}{(81x^4y^{-2})^{14}}$$

\vspace{3cm}
\atsakymas

\question Suprastinkite reiškinį $(1-2x)^2(x+1)-(x+1)^3-3(x+2)(x-1)(x+3)+14x ^{2} -17
$ ir užrašę jį pavidalu $a(x+b)^2+c$ raskite:
\begin{parts}
    \part[7] didžiausią arba mažiausią jo reikšmę bei nurodykite $x$ reikšmę su kuria ta reikšmė gaunama;
   
    \droppoints
    \vspace{10cm}
    \atsakymas
    \newpage
    \part[3] apskaičiuokite $b \cdot \left(a ^{1\overbrace{0\ldots 0}^{666}2} + 21,25 : c\right)    $ .
    \droppoints
    \vspace{2cm}
    \atsakymas
\end{parts}


\question[5] Raskite parametrus $a$, $b$ ir $d$ kad kiekviena lygybė galiotų su visomis $x$ reikšmėmis.
\droppoints
$$ (ax-a)(2x ^{2} -bx+c) = -6 x^3+18 x^2-27 x+15 $$  

\vspace{5cm}
\atsakymas

\question[4] Raskite žemiau esančio reiškinio mažiausią reikšmę:
\droppoints
$$9 x^2+4 y^2+z^2-6 x-12 y+8 z+27$$

\vspace{5cm}
\atsakymas

\end{questions}

\vfill 

\begin{center}
\textbf{Vertinimo lentelė (iš viso 30 taškų)}\par\medskip
\setlength{\tabcolsep}{5pt}
\renewcommand{\arraystretch}{1.3}
\begin{tabular}{|c|*{10}{c|}}
\hline
Taškai & 0--3 & 4--6 & 7--9 & 10--12 & 13--15 & 16--18 & 19--22 & 23--25 & 26--28 & 29--30\\\hline
Pažymys & 1 & 2 & 3 & 4 & 5 & 6 & 7 & 8 & 9 & 10\\\hline
\end{tabular}
\end{center}

\end{document}
